\chapter{Three Complete Projects}
\label{chap:capstones}
\index{capstone project}\index{template!article}\index{template!thesis}\index{template!technical report}

The capstones turn isolated commands into complete scholarly workflows. Each
project begins with a task, records planning decisions, uses one of the supplied
independent source trees, compiles to a PDF, and ends with a release audit. The
sample evidence is simulated; the projects teach structure rather than claim
new scientific results.

\begin{learningoutcomes}
  \item Plan and complete an article, thesis, and technical report.
  \item Navigate each supplied source tree and explain important commands.
  \item Compile all three projects independently from clean folders.
  \item Replace placeholders without inventing institutional or scholarly facts.
  \item Apply document-specific final checklists and extension tasks.
  \item Choose a document form from audience, purpose, evidence, and authority.
\end{learningoutcomes}

\section{How to use the capstones}

Choose the form closest to your immediate work, but inspect all three. They
share equations, tables, figures, citations, and reproducibility practices while
serving different readers:

\begin{center}
\begin{tabularx}{\textwidth}{@{}p{.18\textwidth}XXX@{}}
  \toprule
  Form & Primary audience & Governing constraint & Main outcome \\
  \midrule
  Article & Journal readers and reviewers & Current journal workflow & Bounded scholarly claim \\
  Thesis & Examiners and institution & Degree regulations & Sustained original argument \\
  Report & Named decision owner & Commission and controls & Traceable recommendation \\
  \bottomrule
\end{tabularx}
\end{center}

The template folders contain complete source, not excerpts from this chapter.
Build them before editing. Preserve an untouched copy so that a failed
experiment can be compared with the known-good baseline.

\begin{goodpractice}
For each capstone, keep a short decision log: requirement, choice, reason,
evidence, and date. The log separates an externally mandated rule from a local
preference and makes later migration easier.
\end{goodpractice}

\section{Project 1: A publication-ready scholarly article}

\subsection{The task}

Prepare a generic article titled ``Agreement between two simulated measurement
methods''. It must state that the data are simulated, define the signed paired
difference, report the result in prose and a professional table, include an
original figure, cite authentic sources, and provide complete declarations.
Then create an anonymous review variant and a clean submission rehearsal.

``Publication-ready'' here means structurally and technically ready to adapt.
It does not mean accepted by every journal or scientifically publishable as
empirical research. A real target journal and valid study are still required.

\subsection{Planning notes}

Record before editing:

\begin{itemize}
  \item target article type and intended journal placeholder;
  \item question, unit, definition of $d_i$, and allowed claims;
  \item abstract word limit and section pattern to verify later;
  \item authorship, affiliations, correspondence, and anonymity plan;
  \item required declarations, data objects, figure files, and bibliography;
  \item build command and release evidence.
\end{itemize}

The generic template uses \code{article} because no publisher has been selected.
This decision is explicit and reversible.

\subsection{Expected folder structure}

\begin{terminalcode}{Complete article source tree}
templates/scholarly_article/
|-- main.tex
|-- sections/
|   |-- abstract.tex
|   |-- introduction.tex
|   |-- methods.tex
|   |-- results.tex
|   |-- discussion.tex
|   |-- declarations.tex
|   `-- supplementary_methods.tex
|-- figures/method_comparison.pdf
|-- references.bib
`-- README.md
\end{terminalcode}

\subsection{Complete source and important commands}

The complete editable source is in
\file{templates/scholarly\_article/}. The root file loads packages, registers
\file{references.bib}, defines title metadata, and inputs the section files. The
conditional

\begin{latexcode}{Review identity switch}
\newif\ifanonymous
\anonymousfalse
\end{latexcode}

selects the visible author block. It does not remove identity from all files or
metadata, so the README requires a wider audit.

Methods defines $d_i=B_i-A_i$ and the results refer to table and figure labels
with \cmd{cref}. \pkg{siunitx} typesets quantities; \pkg{booktabs} supplies table
rules; \pkg{biblatex} and Biber produce the reference list. Declarations retain
placeholders where facts cannot be inferred.

\subsection{Compilation instructions}

\begin{terminalcode}{Compile Project 1}
cd templates/scholarly_article
latexmk -pdf main.tex
\end{terminalcode}

Read the console and \file{main.log}; then inspect \file{main.pdf}. A clean build
needs \LaTeX{} and Biber passes, which \pkg{latexmk} schedules.

\subsection{Final article checklist}

\begin{chapterchecklist}
  \item Target journal and current article instructions are recorded.
  \item Title, abstract, results, and conclusion make the same bounded claim.
  \item Author identities and contribution statements are approved.
  \item The anonymous package has been audited beyond the author line.
  \item Data, code, ethics, funding, interest, and acknowledgement statements are accurate.
  \item Table, figure, citations, supplement, units, and labels resolve.
  \item A clean package builds and its portal-style proof has been inspected.
\end{chapterchecklist}

\subsection{Extension exercises}

\begin{enumerate}
  \item Add a second panel that shows signed difference against pair mean.
  \item Generate the summary values from a script rather than fixed source.
  \item Add a reporting-guideline checklist as a named supplement.
  \item Migrate a copy into an official journal class and document every change.
  \item Create separate anonymous and accepted-release manifests.
\end{enumerate}

\section{Project 2: A complete thesis or dissertation}

\subsection{The task}

Prepare a generic five-chapter thesis that uses simulated method-agreement data
to demonstrate a sustained structure. Include title and administrative
placeholders, dedication, acknowledgements, abstract, contents, lists,
abbreviations, literature context, methods, results, conclusion, appendix, and
bibliography. Preserve explicit warnings wherever institutional wording must be
verified.

\subsection{Planning notes}

Build an institutional compliance matrix before changing layout. Identify:

\begin{itemize}
  \item official class, thesis manual, degree, institution, and deposit route;
  \item required order and wording of front matter;
  \item page size, margins, spacing, font, numbering, and accessibility rules;
  \item integrated or thesis-by-publication structure;
  \item joint-work, prior-publication, ethics, copyright, and permissions records;
  \item appendix, data, code, repository, embargo, and file-size requirements.
\end{itemize}

Do not fill a certificate or signature with invented approval. A visible
placeholder is the correct state until authorised information exists.

\subsection{Expected folder structure}

\begin{terminalcode}{Complete thesis source tree}
templates/thesis/
|-- main.tex
|-- metadata.tex
|-- frontmatter/
|-- chapters/
|-- appendices/
|-- figures/
|-- references.bib
`-- README.md
\end{terminalcode}

The full tree contains individual files for all front-matter elements, five
chapters, and supplementary material. \file{metadata.tex} owns repeated identity
fields.

\subsection{Complete source and important commands}

The complete source is in \file{templates/thesis/}. The structural sequence is

\begin{latexcode}{Thesis numbering architecture}
\frontmatter
% preliminary pages and generated lists
\mainmatter
% numbered chapters
\appendix
% lettered appendices
\backmatter
% reference list
\end{latexcode}

\cmd{input} inserts small front-matter units; \cmd{include} gives each chapter
its own auxiliary file and page start. \cmd{listoffigures} and
\cmd{listoftables} read caption data from earlier passes. \cmd{printbibliography}
prints the cited sources and adds them to the contents using its documented
heading option.

The project uses a manual abbreviations table because its list is small. A real
thesis with hundreds of terms may adopt a glossary tool and document its extra
build dependency.

\subsection{Compilation instructions}

\begin{terminalcode}{Compile Project 2}
cd templates/thesis
latexmk -pdf main.tex
\end{terminalcode}

Inspect Roman and Arabic page numbering, title and declaration links, contents,
lists, running heads, appendix letter, citations, and the final reference list.
Perform one build from a fresh copy without \cmd{includeonly}.

\subsection{Final thesis checklist}

\begin{chapterchecklist}
  \item The current institutional class, manual, and deposit rules are authoritative.
  \item Every identity field and required administrative statement is verified.
  \item Front, main, appendix, and back matter follow the mandated order.
  \item Chapter arguments connect and the synthesis addresses the thesis question.
  \item Joint work, prior publication, permissions, and contributions are declared.
  \item Figures, tables, equations, lists, references, data, and code are reproducible.
  \item The complete PDF passes compliance, accessibility, and deposit rehearsal.
\end{chapterchecklist}

\subsection{Extension exercises}

\begin{enumerate}
  \item Replace the manual abbreviation table with a generated glossary workflow.
  \item Add a sixth chapter and a cross-chapter synthesis table.
  \item Implement an institution-supplied title page without duplicating metadata.
  \item Create a thesis-by-publication contribution statement and permission log.
  \item Build a public repository package that safely excludes restricted data.
\end{enumerate}

\section{Project 3: A professional technical report}

\subsection{The task}

Prepare a decision-oriented report on the same simulated measurement comparison.
The report must contain controlled cover metadata, document history,
distribution information, an executive summary, scope, methods, findings, a
table and original figure, limitations, owned recommendations, references, an
evidence register, and a release checklist.

The report must say that the simulated analysis cannot approve an operational
method. Its recommendation is a controlled validation with a prespecified
tolerance.

\subsection{Planning notes}

Identify the commissioning statement, decision owner, primary and technical
audiences, deadline, scope, exclusions, evidence classes, acceptance criteria,
reviewers, approvers, distribution, and archival repository. Define version and
status conventions before preparing the cover.

Separate the analysis finding from a tolerance judgement and the recommendation
from the authorised decision. This distinction is the organising logic of the
report.

\subsection{Expected folder structure}

\begin{terminalcode}{Complete report source tree}
templates/technical_report/
|-- main.tex
|-- sections/
|   |-- document_control.tex
|   |-- executive_summary.tex
|   |-- context.tex
|   |-- methods.tex
|   |-- findings.tex
|   |-- recommendations.tex
|   |-- limitations.tex
|   `-- appendices.tex
|-- figures/difference_plot.pdf
|-- references.bib
`-- README.md
\end{terminalcode}

\subsection{Complete source and important commands}

The complete source is in \file{templates/technical\_report/}. Its
\code{titlepage} creates the controlled cover. Roman numbering covers control
material and lists;

\begin{latexcode}{Report page-number transition}
\pagenumbering{roman}
% control material and contents
\clearpage
\pagenumbering{arabic}
% decision report body
\end{latexcode}

starts the body at page 1. A restrained \pkg{fancyhdr} style repeats project and
classification placeholders. \pkg{tabularx} allows findings and action tables to
use the page width. The action plan connects each recommendation with an owner
and completion evidence.

\subsection{Compilation instructions}

\begin{terminalcode}{Compile Project 3}
cd templates/technical_report
latexmk -pdf main.tex
\end{terminalcode}

Inspect the cover and header classification together, compare the executive
summary with the body, verify table wrapping and figure legibility, and ensure
appendix evidence IDs resolve to preserved objects.

\subsection{Final report checklist}

\begin{chapterchecklist}
  \item The commissioning statement and authorised decision owner are named.
  \item Title, ID, version, date, status, and classification agree.
  \item Executive-summary findings and conditions match the technical body.
  \item Scope, assumptions, criteria, deviations, uncertainty, and limitations are visible.
  \item Recommendations have accepted owners and objective completion evidence.
  \item Technical review, approval, permissions, and distribution are recorded.
  \item The immutable issued PDF, source, evidence, and manifest are archived.
\end{chapterchecklist}

\subsection{Extension exercises}

\begin{enumerate}
  \item Add a risk register with defined likelihood and impact categories.
  \item Create a landscape appendix for a genuinely wide assurance matrix.
  \item Add a requirement-to-evidence traceability table.
  \item Define central commands for report ID, status, version, and issue date.
  \item Produce public and restricted distribution variants without changing findings.
\end{enumerate}

\section{Compare the same evidence across forms}

The simulated mean difference appears in all three projects, but its rhetorical
job changes. The article presents a scholarly result and limitation. The thesis
places it inside a sustained research argument. The report asks what decision
can responsibly follow. This is not merely formatting: audience and authority
change what must be stated.

\begin{scienceinpractice}
One environmental monitoring programme might produce a methods article, a
doctoral thesis containing several site studies, and an agency report recommending
sampling changes. Shared analysis can support all three, yet each document needs
its own scope, provenance, approvals, and release record.
\end{scienceinpractice}

\section{Compile every example from one control point}

The book repository supplies scripts that compile all standalone examples and
all three templates. From the project root:

\begin{terminalcode}{Repository-wide build checks}
./scripts/compile_all_examples.sh
./scripts/compile_templates.sh
./scripts/check_project.sh
\end{terminalcode}

These checks complement, not replace, reading the PDFs. A passing compiler
cannot confirm that an author contribution statement is true, a thesis
declaration is authorised, or a report recommendation has an owner.

\begin{commonerror}
\textbf{Symptom:} a capstone compiles after editing but its claims still refer to
the simulated example while the table now contains real project values.

\textbf{Cause:} content was replaced file by file without a results-consistency
and placeholder audit.

\textbf{Repair:} regenerate evidence from one analysis source, trace each value
and limitation across the document, search all placeholders and simulation
labels, then obtain scientific review.
\end{commonerror}

\section{Code lab: compile all three capstones}

\begin{codelab}
Treat compilation as a test, not a ceremonial final step. Run each build from
its own folder, open every PDF, and record the result beside the relevant
checklist. A capstone is complete only when both source and output survive the
test.
\end{codelab}

\begin{terminalcode}{Project 1: scholarly article}
cd templates/scholarly_article
latexmk -pdf main.tex
\end{terminalcode}

\begin{terminalcode}{Project 2: thesis or dissertation}
cd templates/thesis
latexmk -pdf main.tex
\end{terminalcode}

\begin{terminalcode}{Project 3: technical report}
cd templates/technical_report
latexmk -pdf main.tex
\end{terminalcode}

For a genuine submission, replace the generic build command when a publisher
or institution specifies a different engine or bibliography sequence. The
project README should state that exception explicitly.

\chapterproject

\begin{miniproject}
Build all three templates without editing them. Choose one and copy it into a
new practice folder. Write a one-page project brief, complete its planning
notes with fictional but explicitly labelled values, make one structural and
one evidence change, and run the final checklist. Preserve the baseline and a
diff or change record explaining why every modification was made.
\end{miniproject}

\chaptersummary

The capstones provide complete, independent source for a scholarly article,
thesis, and technical report. Each project connects a task and planning record
to a transparent folder tree, important commands, build procedure, final audit,
and extensions. The same \LaTeX{} mechanisms serve all three, but audience,
governing requirements, and decision authority change their scholarly form.
Successful compilation is the beginning of final verification, not its end.

\chaptercheck
\begin{chapterchecklist}
  \item All three unmodified template baselines compile independently.
  \item The chosen form matches audience, purpose, evidence, and authority.
  \item Planning notes distinguish requirements, conventions, and preferences.
  \item Complete source, assets, bibliography, README, and build command are present.
  \item Simulated evidence and unverified placeholders remain clearly labelled.
  \item Document-specific scientific and release checklists are complete.
  \item The edited capstone can be reconstructed from its archived project.
\end{chapterchecklist}

\chapterexercisestitle
\begin{chapterexercises}
  \item Match six scholarly scenarios to article, thesis, or report and justify each.
  \item Identify the shared and form-specific components of the three projects.
  \item Explain five important commands from one template's root file.
  \item Compile every template and compare page architecture and reference handling.
  \item Replace one placeholder category without creating unverifiable facts.
  \item Design a clean release manifest for each project.
  \item Conduct a peer audit using the appropriate final checklist.
  \item Challenge: derive two document forms from one reproducible analysis while
        preserving their different claims, audiences, and approvals.
\end{chapterexercises}
